This article shows that the use of a single abstract interface to
interact with a blockchain allows one
to write testcases for smart contracts once,
and reuse them against different implementations of the testing blockchain.
Namely, blockchains might consists of a varying number of peers,
with different consensus algorithms. There are local peers (directly contacted
as an object in RAM) and remote peers (contacted through a network).
The same blockchain interface respect the open-closed software engineering
principle. Moreover, by stubbing the abstract blockchain interface,
it becomes possible to execute the
testcases quickly enough to allow their use during continuous integration.
An interesting aspect of using a single abstract interface for different
peer implementations is that testcases become non-regression tests
for modifications to the blockchain, for instance when replacing its
consensus algorithm, as it actually happened during the development of the
Hotmoka blockchain.
